% =====================================================================
%  Sec. 7 -- Three prices: support, resolution, and shots
%  Replaces the subsection "Scaling of the sampling cost with system size"
%  of the former paper/results.tex (lines 251-299 of the v2 tree).
%
%  NEW \label's introduced here (none exist in the v2 tree):
%      sec:prices, eq:supportexp, tab:fracgrid, tab:shots, tab:finiteshot
%  \cite keys: none new are needed in this section.
%  External labels used: sec:leakage, sec:moments, conj:geom, app:moments:conj,
%      sec:frontier, sec:method, fig:scaling.
%      *** \ref{thm:b2} has been REMOVED from this file: Sec. 4 withdraws the
%      corollary that sentence invoked, and theorem_b2.tex is superseded. ***
%  FLOAT rebuilt in parallel, NOT written here:
%      \ref{fig:scaling} -- the unlabelled shaded band and its inked label at
%      fig_scaling2_native.tex:21,30 removed; finite-shot band with +-sigma
%      added; secondary shot axis; |S| = 19183 / 824503.
%
%  PROVENANCE OF THE FINITE-SHOT NUMBERS (Sec. 7.3 and 7.4).  TWO runs, and they
%  are named as two throughout:
%    run A  release/data/sampled_honest.json  (sampled_honest.py, 2026-09-18T11:57Z,
%           8 seeds 20260918..20260925, L = 4, 6, 8 only, K=18, dt=0.5, eta=0.15t,
%           n_l=260, 600-point grid).
%    run B  release/data/honest_sampling.json (honest_sampling_scaling.py,
%           2026-09-18T19:07Z, seeds 7001.., L = 4..12, 4 seeds per size and 6 at
%           L=10, same protocol, nl_exact=260 / nl_sub=150).
%  BOTH ARE NOW DEPOSITED, with their generating scripts in release/src/; the
%  %TODO(deposit) that stood here is discharged.  The earlier version of this
%  section said that nothing was measured above L=8.  That was wrong: run B
%  measures L=10 and L=12, and its L=6 and L=8 budgets agree with run A's to
%  0.2% and 3.2%.  Table tab:shots now carries five measured rows.
% =====================================================================

\section{Three prices: support, resolution, and shots}\label{sec:prices}

On half-filled periodic Hubbard rings at $U/t=8$ ($L=4$--$14$, $8$ to $28$ qubits) the method is
priced on three axes, each measured here: the determinant support the sampler must populate grows
as $e^{1.05\,L}$, more slowly than the sector's $e^{1.30\,L}$; the sector fraction, unlike the
support, is a property of the question put to the model rather than of the model, its exponent
moving by a factor $2.2$ over a decade of the resolution $\eta$; and the shot budget $T$ that
populates the support, measured at $L=4$--$12$, lies two orders of magnitude above it and grows as
$e^{1.12\,L}$ over $L=6$--$12$, its $L=14$ value being a floor of $1.3\times10^{8}$ shots for a
single seed and branch.

\paragraph{The dynamical support.} $|\mathcal S|_\eta$ here is the number of
determinants needed for the reconstruction of $A(\omega)$ to reach relative $L_1<\mathrm{thr}$ at
resolution $\eta$ and recursion depth $n_\ell$, so it depends on all three
(\texttt{scaling\_lanczos\_mf.py}, \texttt{stage\_frac}: $\mathrm{thr}=0.05$, $\eta=0.15\,t$,
$n_\ell=150$). It is distinct from the static ground-state support $|\mathcal S|_{\varepsilon_w}$ of
Sec.~\ref{sec:nogo} and from the returned set $S$ of Theorem~\ref{thm:leakage}: the three coincide only
by accident, and no statement here transfers to the other two (Sec.~\ref{sec:intro}).

\subsection{The support exponent and its resolution dependence}

Across $L=4$--$14$ ($8$ to $28$ qubits, half-filled \emph{periodic} rings at $U/t=8$; the three largest
sizes reached by matrix-free Lanczos/Haydock and validated against dense diagonalization at $L\le8$),
the determinant support needed to reconstruct $A(\omega)$ to relative $L_1<0.05$ at $\eta=0.15\,t$ and
$n_\ell=150$ grows exponentially,
\begin{widetext}
\begin{equation}
|\mathcal S|_{\eta=0.15} \;=\; 22,\;246,\;2180,\;19183,\;124407,\;824503
\;\;\Longrightarrow\;\; |\mathcal S|_\eta\sim e^{1.05\,L}\quad(L=4\text{--}14,\ R^{2}=0.998),
\label{eq:supportexp}
\end{equation}
\end{widetext}
more slowly than the $(N{+}1)$ sector itself,
$\dim = 24,\,300,\,3920,\,52920,\,731808,\,10306296\sim e^{1.30\,L}$. Restricted to $L=4$--$12$, the five sizes of the replication grid, the fit gives $1.082$ at
$R^{2}=0.998$. The fraction is the support divided by the sector dimension, exactly, so the fraction and
support exponents differ by the constant $b_{\dim}=1.291$ and move together, by $0.148$ over the decade
$\eta=0.05\to0.50\,t$ (Table~\ref{tab:fracgrid}): a seventh of the support exponent, three quarters of
the fraction exponent. We therefore report the support, the absolute resource a device has to populate,
and the fraction as the derived, resolution-dependent quantity that it is.

\subsection{The sector fraction is a property of the question, not of the model}

At $\eta=0.15\,t$ the fraction needed for a fixed spectral accuracy falls from $0.92$ at $L=4$ to
$0.08$ at $L=14$, but its exponent is set by $\eta$: over the decade of $\eta$ it moves by a factor
$2.2$, from $-0.1266$ to $-0.2743$ (Table~\ref{tab:fracgrid}), and even at fixed $\eta$ the local
logarithmic slopes between consecutive published points steepen by a factor $6.8$ across five
intervals, so the fraction is not a single exponential either. The reference is itself a
depth-$150$ Lanczos computation whose drift between depths $100$ and $200$ at $L=12$ is $0.93\%$ at
$\eta=0.15\,t$ and $3.0\%$ at $\eta=0.10\,t$, so the fraction is reported as a declared window,
$\eta\ge0.15\,t$ at $n_\ell=150$, with the full grid and the finite-depth caveats in
Sec.~\ref{sm:prices}.

\subsection{The third price: shots}

Populating a subspace by measurement is a coupon-collector problem on the Born distribution of the
time-evolved seed, so the budget exceeds the support it buys. Table~\ref{tab:shots} reports that budget
\emph{measured}, by multinomial sampling of the declared protocol, at $L=4$--$12$, in
two independent runs that agree where they overlap to $0.2\%$ and $3.2\%$; every budget is for
\emph{one} orbital seed and \emph{one} branch.

\begin{table*}[tp]\centering\small
\caption{\textbf{The third price, measured at every size but the largest.} $T$ is the total shot budget,
split evenly over the $19$ snapshots of the protocol, for a single orbital seed and a single branch; it
is the budget at which the finite-shot reconstruction itself first reaches relative $L_1<0.05$, and
$\pm$ is the sample standard deviation across seeds, not the standard error of the mean. Run~A is
\texttt{sampled\_honest.json} ($8$ seeds, $L\le8$); run~B is \texttt{honest\_sampling.json} ($4$ seeds,
$6$ at $L=10$), and the two are separate runs with different seeds and different budget grids, listed
separately at $L=6$ and $L=8$. At $L=4$ run~B
is censored by its own cap of $10^{8}$ shots per snapshot, the $24$-state sector being degenerate for
this purpose, and run~A is quoted there. The last column divides $T$ by the \emph{published} support;
dividing instead by the support each run actually populated gives $95$, $107$, $141$ and $124$ at
$L=6$--$12$. Only $L=14$ is a floor: the overhead held at its $L=12$ value, which is a lower bound and
not an estimate.}\label{tab:shots}
\setlength{\tabcolsep}{5pt}
\begin{tabular}{lccccl}
\toprule
$L$ & $\dim(N{+}1)$ & $|\mathcal S|_\eta$ & $T$ & $T/|\mathcal S|_\eta$ & status\\
\midrule
$4$  & $24$           & $22$       & $(7.8\pm2.0)\times10^{2}$   & $35$  & measured, run~A, $8$ seeds\\
$6$  & $300$          & $246$      & $(2.43\pm0.34)\times10^{4}$ & $99$  & measured, run~A, $8$ seeds\\
$6$  & $300$          & $246$      & $(2.44\pm0.23)\times10^{4}$ & $99$  & measured, run~B, $4$ seeds\\
$8$  & $3\,920$       & $2\,180$   & $(2.70\pm0.13)\times10^{5}$ & $124$ & measured, run~A, $8$ seeds\\
$8$  & $3\,920$       & $2\,180$   & $(2.61\pm0.08)\times10^{5}$ & $120$ & measured, run~B, $4$ seeds\\
$10$ & $52\,920$      & $19\,183$  & $(3.09\pm0.04)\times10^{6}$ & $161$ & measured, run~B, $6$ seeds\\
$12$ & $731\,808$     & $124\,407$ & $(1.91\pm0.01)\times10^{7}$ & $154$ & measured, run~B, $4$ seeds\\
$14$ & $10\,306\,296$ & $824\,503$ & $\gtrsim1.3\times10^{8}$    & $154$ & floor, not measured\\
\bottomrule
\end{tabular}
\end{table*}

The budget is two orders of magnitude larger than the support it buys, already at $L=6$, and it
grows exponentially, $7.8\times10^{2}$ to $1.9\times10^{7}$ over $L=4$--$12$, a log-linear fit over
$L=6$--$12$ giving $T\sim e^{1.12\,L}$ at $R^{2}=0.996$. At $L=8$ the sixteen-channel momentum-resolved
configuration of Fig.~\ref{fig:akwsampled} reached its own resolution and
accuracy with $8.3\times10^{7}$, and not with half as many; the $16\times2.7\times10^{5}\approx4\times10^{6}$
shots that sixteen channels at the local-seed budget would cost is a transfer and not a measurement,
because the momentum seeds have larger support. The $L=14$ floor of $1.3\times10^{8}$ shots for a
\emph{single} seed and branch (Table~\ref{tab:shots}) exceeds the largest shot budget drawn anywhere
in this work, the $8.3\times10^{7}$ shots of the sixteen-channel configuration of
Fig.~\ref{fig:akwsampled}. The overhead per determinant is a level and not a growth law---$35$,
$99$, $124$, $161$, $154$ against the published support across $L=4$--$12$, rising by a factor $3.5$
over the first three sizes ($35$ to $124$), then to $161$ at $L=10$, and not rising over the last
step, $161$ to $154$. Its exponent is not resolved by these data. Fitted on $L=6$--$12$ it is $0.08$ per site against the published support ($95\%$ interval $[-0.03,0.19]$) and $0.054$ ($[-0.066,0.174]$) against the support each run populated; both intervals contain zero, so \emph{these data do not establish that the shot budget grows with a larger exponent than the
support}, only that it is larger in absolute terms at every size measured.

\paragraph{What the published curve is.} The published curve is the $T\to\infty$ limit of the
declared protocol: the ranking that produced the published fractions is \texttt{argsort} applied to
the Born distribution of the time-evolved state accumulated over the $19$ snapshots of the
protocol---the distribution the device would itself measure with infinitely many shots---and not a
clairvoyant ordering by exact eigenvector amplitudes. With
finite shots the fraction needed rises, from $0.8200$ to $0.8571\pm0.0135$ at $L=6$ and from $0.5561$ to
$0.6275\pm0.0031$ at $L=8$ (Table~\ref{tab:finiteshot}), so the published curve is a strict lower bound
on the finite-shot cost. At fixed size the penalty is budgetary, not structural: at $L=8$ raising the budget by a
factor $12.7$ takes the ratio of the finite-shot error to the infinite-shot-ordering error from
$1.96\pm0.12$ to $1.06\pm0.01$. Finite-shot fractions are therefore reported only with their budget and
seed count attached.

